\documentclass[11pt,a4paper]{article}
\usepackage[T1]{fontenc}
\usepackage{lmodern,amsmath,amssymb,booktabs,microtype}
\usepackage[margin=24mm]{geometry}
\usepackage[colorlinks=true,urlcolor=blue,linkcolor=blue]{hyperref}
\newcommand{\EE}{\mathbb E}
\newcommand{\tr}{\operatorname{tr}}
\newcommand{\Rem}{\operatorname{Rem}}
\newcommand{\W}[1]{\mathsf W_{#1}}
\newcommand{\core}{\mathcal E}
\title{Restoring the raw memory descendants\\A two-chaos residual for the second memory coefficient}
\author{Alan Day \and Chengheng\thanks{AI research assistant (OpenAI).}}
\date{Local research note --- 6 November 2026}
\begin{document}\maketitle
\begin{abstract}
We restore the lower Gaussian-chaos components of the four previously
retained quartic descendants. Their second-chaos limits contain no
output-deficit direction. A fixed-charge trace argument transfers
the earlier fourth-chaos Gram to half filling. Consequently the raw
second memory coefficient, after projection onto a specified finite
span, retains both the known fourth-chaos residual and the newly
resolved second-chaos residual. This completes their equal-time
assembly without assuming that their energy-resolved cross vanishes.
\end{abstract}
\section{State, scale, and result}
Use standard CAR and the full overlapping complex Gaussian $p$-body
Hamiltonian, $\hbar=1$, fixed charge $m=N/2$, normalized maximally
mixed trace, and $\langle X,Y\rangle=\EE\tr_m X^\dagger Y$.
All projections use this annealed inner product. Write $h=H_m$,
$k=H_{m-1}$ with the same couplings, $R_j=\Rem(k^j)$, and
\begin{equation}
 A=R_2+(2p-m)h^2,\quad B=\W2h^2,\quad O=\W2A,
 \qquad J_hX=\{h,X\}/2.
\end{equation}
Take $\lambda_N=p^2/N\to\lambda>0$, $q_N=e^{-4\lambda_N}$,
$C_N=\EE\tr_mh^2\to\sigma^2=J_s^2/\lambda$. Lower-chaos
coefficients retain $C_N,\lambda_N,q_N$ until after their growing
terms are subtracted. Let
\begin{equation}
 \mathcal V_N=\operatorname{span}
 \{I,B,O,J_hF,E_0,M_0,S_0,Z_0\},
 \qquad \widehat\Pi_N=\operatorname{proj}_{\mathcal V_N},
\end{equation}
where the four raw quartic descendants are defined on the next page.
With $\Pi_N=\operatorname{proj}_{\{I,B,O,J_hF\}}$ and
$T_N=(1-\Pi_N)\mathcal S_N$, the result is
\begin{equation}
 (1-\widehat\Pi_N)mZ_2
 =Q_{4,N}+8\lambda_N C_Nq_N^2T_N+o_{L^2}(1).\label{main}
\end{equation}
$Q_{4,N}$ is the preceding fourth-chaos residual and
$\mathcal S_N$ is the second-chaos output-deficit insertion.
Their equal-time inner product is exactly zero by Gaussian-chaos
orthogonality. Thus the limiting norm squared is $V_4+V_2>0$,
with explicit expressions below. The claim concerns this operator
span, not the full algebra of conserved quantities.

\newpage
\section{Five raw words and exact low-chaos restoration}
Define $\mathcal T_j=R_j+(jp-m)h^j$. The five raw operators are
\begin{align}
 E_0&=\{A,h^2\}/2,&M_0&=hAh,\notag\\
 S_0&=\{h,\mathcal T_3\}/2-E_0-M_0,
 &C_0&=\mathcal T_4-2E_0-M_0-2S_0,\\
 Z_0&=h^4.\notag
\end{align}
They are linear combinations of
$\mathbf R=(R_4,\{h,R_3\},\{h^2,R_2\},hR_2h,h^4)^T$:
\begin{equation}
 \begin{pmatrix}E_0\\M_0\\S_0\\C_0\\Z_0\end{pmatrix}
 =\begin{pmatrix}
 0&0&1/2&0&2p-m\\
 0&0&0&1&2p-m\\
 0&1/2&-1/2&-1&m-p\\
 1&-1&0&1&0\\
 0&0&0&0&1
 \end{pmatrix}\mathbf R.\label{matrix}
\end{equation}
In particular, the memory core satisfies the exact identity
\begin{equation}
 \core=C_0-2S_0+(E_0+M_0)/2.
\end{equation}
The single-contraction map for a homogeneous quartic polynomial is
$\mathfrak C=\frac12\sum_i v_i\partial_{x_i}^2$, where $x_i$ are
independent real coupling coordinates of variance $v_i$. Hence
\begin{equation}
 e_P=\EE P,\qquad \W2P=\mathfrak CP-2e_P,
 \qquad \W4P=P-\mathfrak CP+e_P.\label{chaos}
\end{equation}
Apply (\ref{matrix}) to the preceding five exact contractions
$(D_4,D_{13},D_{22},D_{121},D_{40})^T$ to obtain every
$\mathfrak CP$. This uses the same-coupling adjacent-charge
covariance, not independent Hamiltonians \cite{kota,local}.

For clarity, each scalar mean can also be recovered from the rank
filters. If $\mathfrak CP=\sum(g_P^+Z^++g_P^-Z^-)$, then in the
common-rank domain $p\leq\min(m-1,N-m)$,
\begin{equation}
 e_P=\frac12\sum_{a=0}^p
 \left[C_a g_P^+(a,a,0)+mC_a^-g_P^-(a,a,0)\right]I.\label{mean}
\end{equation}
Here $C_a=\EE\tr_m h_a^2$ and $C_a^-=\EE\tr_{m-1}k_a^2$.
The common-rank condition holds eventually in the present limit.
Equations (\ref{matrix})--(\ref{mean}) restore $\W0,\W2,\W4$
at finite $N$; scalar subtraction is not left implicit.

\newpage
\section{The complete second-chaos descendant table}
For $P=E_0,M_0,S_0,C_0,Z_0$, write
\begin{equation}
 \W2P=C_N[p\ell_P+k_P]B+C_No_PO+j_PJ_hF
                         +C_Ns_P\mathcal S_N+o_{L^2}(1).\label{tableformula}
\end{equation}
The table uses $\lambda=\lambda_N$ and $q=q_N$ on its right-hand
side, so the coefficients of order $p$ retain finite-$N$ parameters.
\begin{center}\small
\begin{tabular}{cclccc}\toprule
 $P$&$\ell_P$&$k_P$&$o_P$&$j_P$&$s_P$\\\midrule
 $E_0$&$2$&$-2\lambda q^2$&$1$&$0$&$0$\\
 $M_0$&$2$&$-2\lambda(2q-1)$&$q^2$&$0$&$0$\\
 $S_0$&$2q$&$-2\lambda(q^2-q+1)$&$q$&$1$&$0$\\
 $C_0$&$2q^2$&$-2\lambda(8\lambda q^2-3q^2+2q+2)$&$1$&$2q$&$8\lambda q^2$\\
 $Z_0$&$0$&$q^2+2q+3$&$0$&$0$&$0$\\\bottomrule
\end{tabular}
\end{center}
This is obtained by applying (\ref{matrix}) to the exact filters,
then changing from $(Z^+,Z^-)$ to $(Z^+,A_{ab;r})$ before expanding.
The full rational coefficients and both fixed-deficit and
rank-proportion expansions are saved in the symbolic certificate.

Here is the limit justification for all five rows. Each $A$-filter
has bounded rational prefactors multiplying
$1,\kappa_a/C_N,\kappa_b/C_N,\kappa_r/C_N$.
Their adjacent-rank difference is $O(p^{-1})$, giving
\begin{equation}
 |g_{A,P}/C_N-o_P|\leq K_\lambda(1+u+v+S)/p,
 \quad |g_{A,P}/C_N|\leq K_\lambda.
\end{equation}
The preceding input-weighted $A$ norm and sixth ordinary-product
moment bound therefore apply unchanged: a window $S\leq p^{3/4}$
gives an $O_{L^2}(p^{-1/4})$ error for this part.

Each $H$-filter, after subtracting its leading $p\ell_P$, has a
uniform bound $K_\lambda(1+u+v+S)$ and a fixed-deficit affine limit.
The previous factorial input tail and third output moment give
uniform integrability of its square. Thus the $H$-part converges
in $L^2$. Finally,
$\W2\{h,D_N\}=L_{\rm in}+2\bar s_NB$ and
$J_hF-4\lambda_Nq_NC_NL_{\rm in}\to0$ give the centered coefficients
$k_P,j_P$. The total remainder in (\ref{tableformula}) is only
claimed to be $o_{L^2}(1)$.

The decisive entry is the last column. Every old raw descendant
has $s_P=0$; restoring them cannot absorb the output-deficit
direction through an additional second-chaos component.

\newpage
\section{Transferring Wick-block moments to half filling}
The earlier marked-word calculation used $2^{-N}\operatorname{Tr}$.
A bridge is required before using its fourth-chaos Gram here.
For a diagonal occupation constraint on $L$ distinct sites, with
$t$ occupied and $L-t$ empty, the ratio of fixed half-filled trace
to full-Fock trace is exactly
\begin{equation}
 \mathcal R_N(L,t)=
 \frac{2^L(m)_t^{\downarrow}(m)_{L-t}^{\downarrow}}
 {(N)_L^{\downarrow}}
 =2^L\frac{\binom{N-L}{m-t}}{\binom Nm}.\label{ratio}
\end{equation}
Uniformly for $L\leq kp=O(\sqrt N)$ and allowed $t$,
\begin{equation}
 \log\mathcal R_N(L,t)
 =\frac{L-(2t-L)^2}{2N}+O(L^3/N^2),\qquad
 \mathcal R_N(L,t)\leq \exp(K_kp^2/N).\label{ratioexpand}
\end{equation}
The expansion follows from the three falling factorials; the upper
bound also follows directly by dropping their numerator factors
after extracting powers of $m$. It is valid for large $N$ at fixed
word length $k$.

In a Gaussian-paired word, every chord carries two independent
$p$-element index sets. At a site belonging to only one such set,
the paired creation/annihilation entries force its initial occupation.
The two sets of each chord have opposite singleton occupations and
equal cardinality. If $K$ is the total pairwise intersection count,
deleting the multiply used sites changes their balance by at most
a constant times $K$. Therefore $|2t-L|\leq c_kK$ for every
nonzero local occupation term. The unmarked word has $L\leq kp$.

For the marked Wick blocks used here, no Gaussian chord is internal
to a block. Each body-rank mark hence contracts indices from
different chords. It is supported on their intersections and leaves
the singleton argument intact. Expanding the multiply used sites
in occupation projectors has absolute coefficients bounded by
$e^{c_kK}$ times a fixed polynomial in $K$. This is also the
absolute overlap bound used in the preceding marked-word derivation.

First restrict $K\leq K_0$. Equations (\ref{ratio})--(\ref{ratioexpand})
give $\mathcal R_N\to1$ uniformly. The established exponential
moment bounds for intersections of a fixed number of random
$p$-element sets control the complementary absolute tail, including
the marks; the global bound in (\ref{ratioexpand}) preserves it.
Letting $K_0\to\infty$ proves equality of the two trace limits.

This transfers the fourth-chaos Gram and its fixed energy moments.
It also transfers the bounded fixed-degree marked Wick-block
moments needed for the centered fields below. It does not identify
ordinary coupling chaos with the spectral $q$-Wick polynomials of
\cite{tt}; those are different constructions.

\newpage
\section{Projection after restoring both chaos sectors}
Let $D_i=(1-\Pi_N)P_i$ for $P_i=E_0,M_0,S_0,Z_0$.
Their means are scalar and their second-chaos parts lie in
$\operatorname{span}\{B,O,J_hF\}$ up to $o_{L^2}(1)$. Thus
\begin{equation}
 D_i=\W4P_i+o_{L^2}(1).
\end{equation}
Their limiting Gram is the previously computed fourth-chaos
nuisance Gram, now justified in the same half-filled trace.
It is strictly positive. With Gram entries divided by $\sigma^8$,
its determinant is
\begin{align}
 16\lambda^3(1-q)^2(1+q)^4(1+q^2)d_3d_5
 (1+q+4\lambda q)^2\notag\\
 {}\times[1-q^3+4\lambda q^2(2q^2-q+2)]>0,\quad
 d_3=q^2-q+1,\ d_5=1+q+q^2+q^3+q^4.\label{det}
\end{align}
The fourth-chaos residual is
\begin{align}
 Q_{4,N}&=\W4(C_0-\alpha E_0-\beta M_0-\gamma S_0-\delta Z_0),\\
 \alpha&=-\frac{2q^3}{d_3d_5},\quad
 \beta=-\frac{q^2(1+q+q^2)}{d_3d_5},\quad\gamma=\frac{2q}{d_3},\notag\\
 \delta&=\frac{2\lambda(1-q)^3(1+q)(1+q+q^2)}{(1+q^2)d_5}.\notag
\end{align}
These coefficients can be evaluated at $\lambda_N,q_N$.
After applying $1-\Pi_N$, the core is bounded and equals its
fourth-chaos component plus $8\lambda_N C_Nq_N^2T_N+o_{L^2}(1)$.
The latter component is exactly orthogonal to every $\W4P_i$.
Invertibility of the Gram and continuity of this finite-dimensional
projection prove (\ref{main}) for $\core$ in place of $mZ_2$.

The two limiting contributions are
\begin{align}
 V_4&=\frac{J_s^8}{\lambda^4}
 \frac{4\lambda(1-q)^2(1+q)^2[P(q)+4\lambda q^3U(q)]}
 {(1+q^2)d_3d_5},\label{V4}\\
 P(q)&=1+q+2q^2+q^3+q^4,\quad
 U(q)=2+q+2q^2+q^3+2q^4,\notag\\
 V_2&=\frac{64J_s^8q^4(1+q)(1-q^2+8\lambda q)}
 {\lambda(1+q+4\lambda q)}.\label{V2}
\end{align}
Both are positive for fixed $\lambda>0$. After the double-scaled
limit, $V_4\sim768J_s^8/(5\lambda)$ and $V_2\to1024J_s^8$.
These are energy-to-the-eighth variances of the collective
coefficient $mZ_2$, not of unscaled $Z_2$.

\newpage
\section{From the core to the raw second memory coefficient}
For completeness, the second centered memory coefficient is defined
from $\widehat M(s)=G+isZ_1-s^2Z_2/2+O(s^3)$ at finite $N$.
The exact first-return calculation gives, with $r=p/m$,
\begin{equation}
 mZ_2=B_{22}^f+\frac{(1+r)^2}{4}[h,[h,A]]
       +2rJ_hT_3+3r^2J_h^2A-A^2/m-mr^4h^4.
\end{equation}
Here $T_3$ is the Hermitian cubic derivative singlet and
$B_{22}^f=\sum_i(\mathcal L^2c_i)^\dagger\mathcal L^2c_i$.
Use the previous centered fields
\begin{equation}
 K_A=a_m+c_m(h^2-C_N),\quad Y=A-K_A,\quad
 P_m=t_{3,m}h^3+t_{1,m}h,\quad U=T_3-P_m.
\end{equation}
$c_m$ is the exact energy-quadratic regression coefficient; $P_m$
is the preceding cubic conserved subtraction. Its precise scalar
coefficients do not enter the projection below. Define
\begin{equation}
 D_{2,m}=2rJ_hP_m+3r^2J_h^2K_A-K_A^2/m-mr^4h^4.
\end{equation}
It belongs exactly to $\operatorname{span}\{I,h^2,h^4\}\subset
\mathcal V_N$. The remaining identity is
\begin{align}
 mZ_2-D_{2,m}-\core={}&(r/2+r^2/4)[h,[h,Y]]
                       +2rJ_hU+3r^2J_h^2Y\notag\\
 &-\bigl(\{K_A,Y\}+Y^2\bigr)/m.\label{remainder}
\end{align}
The centered Wick-block representations of $Y,U$ transfer by the
argument on page 4. Their linear filtered pieces are also bounded:
for an invariant Hermitian Gaussian field $X$ with $\EE X^2=vI$,
Wick pairing and Cauchy--Schwarz give $\EE\tr X^4\leq3v^2$;
their $v$ is controlled by the previous input-rank tails. The
cubic Wick blocks and the products with $Y$ have the preceding
absolute marked-word bounds. Thus the fixed-sector moments needed
for $h^2Y,hU,Y^2$ are bounded. Noncommutative H\"older uses, for
example, $\EE\tr h^8$, $\EE\tr Y^4$, and $\EE\tr U^4$;
unweighted $L^2$ bounds alone would not suffice.

Since $r=O(p^{-1})$, $a_m=O(p)$, $c_m=O(1)$, and $m=O(p^2)$,
every term on the right of (\ref{remainder}) is $O_{L^2}(p^{-1})$
or smaller. In particular $a_mY/m=O_{L^2}(p^{-1})$;
the potentially finite conserved contribution of $A^2/m$ was
already included in $D_{2,m}$ and has not been dropped.
Projection is contractive and annihilates $D_{2,m}$, completing
(\ref{main}) for the raw coefficient.

Ordinary chaos is orthogonal only for the equal-time annealed
inner product. Inserting $e^{ith}$ or powers of $h$ changes the
Gaussian degree bookkeeping. Thus the $Q_4$ autocorrelation known
from the earlier work does not by itself determine the full residual
spectrum: the $T_N$ autocorrelation and its energy-resolved cross
with $Q_{4,N}$ must still be calculated.

\newpage
\section{Checks and the next spectral question}
The symbolic record contains the exact five-word matrix, all
rational filters, their full rank-proportion expansions, and the
completed second-chaos table. Independent even polarization of
the raw quartic definitions verifies 3,600 contractions in 720
directions. The first two systems, $(N,p,m)=(3,1,2),(4,2,3)$,
cover every quadratic Gaussian coefficient; three additional
directions use the larger half-filled system $(6,2,3)$.
All entries are checked dyadic integers within the exact arithmetic
range before comparison. A further 720 identities check the raw
core decomposition.

Second contraction gives the five scalar means. Their independent
rank sums agree in 10 cases. The $(4,2,3)$ sector has a lower-only
rank absent from the upper sector, so it is outside the common-rank
filter domain; the attempted mean check exposed this and was
recorded. Its direct CAR contractions remain valid. The eventual
half-filled limit lies inside the common-rank domain.

The assembly certificate also checks 123 exact charge-trace
identities, the leading logarithm in (\ref{ratioexpand}), the
fourth-chaos projection equations, the determinant (\ref{det}), and
the two small-$\lambda$ limits. The trace-limit proof uses the
balanced-singleton and absolute-tail argument, not these finite
identity checks alone. An initial symbolic parsing failure involving
the name \texttt{lambda} is retained alongside the corrected run.

The next task is the fixed-half-filled joint two-energy measure of
$T_N$ and $Q_{4,N}$. This will decide how the two positive equal-time
pieces combine in the actual memory spectrum. A central-charge
extension of the second-chaos filter and the finite-memory remainder
are subsequent steps toward the bulk singlet/meson dictionary.

\begin{thebibliography}{9}
\bibitem{kota} V. K. B. Kota and M. Vyas,
\emph{Random Matrix Theory for Transition Strength Densities in Finite
Quantum Systems: Results from Embedded Unitary Ensembles},
\href{https://arxiv.org/abs/1504.00759}{arXiv:1504.00759}, Sec.~2.
\bibitem{tt} J. Tall and S. Tomsovic,
\emph{Double-scaled bosonic and fermionic embedded ensembles,
complex SYK, and the dual Hilbert space},
\href{https://arxiv.org/html/2604.14522v2}{arXiv:2604.14522v2},
Secs.~2.1 and 3.1. The spectral Wick product is defined in Eq.~(3.2).
\bibitem{local} Alan Day and Chengheng, local project notes:
\emph{Second memory coefficient and quartic singlet};
\emph{Quadratic contractions of the memory core};
\emph{Two-charge memory filters}; \emph{Fixed-charge residual norm};
\emph{The quadratic memory core in the double-scaled limit}
(6 November 2026). Source paths are listed in the accompanying
literature record.
\end{thebibliography}
\end{document}
